\begin{lemma}
\label{lemma-regular-immersion-ext}
Let $Y$ be a ringed space. Let $\mathcal{I} \subset \mathcal{O}_Y$
be a sheaf of ideals. Set $\mathcal{O}_X = \mathcal{O}_Y/\mathcal{I}$ and
$\mathcal{N} =
\SheafHom_{\mathcal{O}_Y}(\mathcal{I}/\mathcal{I}^2, \mathcal{O}_X)$.
If $\mathcal{I}$ is Koszul-regular
(Divisors, Definition \ref{divisors-definition-regular-ideal-sheaf}) then
$$
R\SheafHom_{\mathcal{O}_Y}(\mathcal{O}_X, \mathcal{O}_Y) =
\wedge^r \mathcal{N}[-r]
$$
where $r : Y \to \{1, 2, 3, \ldots \}$ sends $y$ to
the minimal number of generators of $\mathcal{I}$ needed in a neighbourhood
of $y$.
\end{lemma}

\begin{proof}
We can use Lemmas \ref{lemma-ext} and \ref{lemma-regular-ideal-ext}
to see that we have isomorphisms
$\wedge^i\mathcal{N} \to
\SheafExt^i_{\mathcal{O}_Y}(\mathcal{O}_X, \mathcal{O}_X)$
for $i \geq 0$. Thus it suffices to show that the map
$\mathcal{O}_Y \to \mathcal{O}_X$ induces an isomorphism
$$
\SheafExt^r_{\mathcal{O}_Y}(\mathcal{O}_X, \mathcal{O}_Y)
\longrightarrow
\SheafExt^r_{\mathcal{O}_Y}(\mathcal{O}_X, \mathcal{O}_X)
$$
and that
$\SheafExt^i_{\mathcal{O}_Y}(\mathcal{O}_X, \mathcal{O}_Y)$
is zero for $i \not = r$. These statements are local on $Y$. Thus
we may assume that we have global sections $f_1, \ldots, f_r$ of
$\mathcal{O}_Y$ which generate $\mathcal{I}$ and which form a
Koszul regular sequence. Let $\mathcal{A}^\bullet$
be the Koszul complex on $f_1, \ldots, f_r$ as introduced in the proof of
Lemma \ref{lemma-regular-ideal-ext}. Then
$$
R\SheafHom_{\mathcal{O}_Y}(\mathcal{O}_X, \mathcal{O}_Y) =
\SheafHom^\bullet(\mathcal{A}^\bullet, \mathcal{O}_Y)
$$
by Cohomology, Lemma \ref{cohomology-lemma-Rhom-strictly-perfect}.
Denote $1 : \mathcal{A}^\bullet \to \mathcal{O}_Y$ the map of differential
graded $\mathcal{O}_Y$-algebras given by the identity map of
$\mathcal{A}^0 = \mathcal{O}_Y \to \mathcal{O}_Y$ in degree $0$.
With $\delta_j$ as in the proof of Lemma \ref{lemma-regular-ideal-ext}
we get an isomorphism of graded $\mathcal{O}_Y$-modules
$$
\mathcal{O}_Y\langle \delta_1, \ldots, \delta_r\rangle
\longrightarrow
\SheafHom^\bullet(\mathcal{A}^\bullet, \mathcal{O}_Y)
$$
by mapping $\delta_{j_1} \ldots \delta_{j_i}$ to
$1 \circ \delta_{j_1} \circ \ldots \circ \delta_{j_i}$ in degree $i$.
Via this isomorphism the differential on the right hand side
induces a differential $\text{d}$ on the left hand side.
By our sign rules we have $\text{d}(1) = - \sum f_j \delta_j$.
Since $\delta_j : \mathcal{A}^\bullet \to \mathcal{A}^\bullet[1]$
is a morphism of complexes, it follows that
$$
\text{d}(\delta_{j_1} \ldots \delta_{j_i}) =
(- \sum f_j \delta_j )\delta_{j_1} \ldots \delta_{j_i}
$$
Observe that we have $\text{d} = \sum f_j \delta_j$ on the differential
graded algebra $\mathcal{A}$. Therefore the map defined by the rule
$$
1 \circ \delta_{j_1} \ldots \delta_{j_i} \longmapsto
(\delta_{j_1} \circ \ldots \circ \delta_{j_i})(\xi_1 \ldots \xi_r)
$$
will define an isomorphism of complexes
$$
\SheafHom^\bullet(\mathcal{A}^\bullet, \mathcal{O}_Y)
\longrightarrow \mathcal{A}^\bullet[-r]
$$
if $r$ is odd and commuting with differentials up to sign if $r$ is even.
In any case these complexes have isomorphic cohomology, which shows the
desired vanishing. The isomorphism on cohomology in degree $r$
under the map
$$
\SheafHom^\bullet(\mathcal{A}^\bullet, \mathcal{O}_Y)
\longrightarrow
\SheafHom^\bullet(\mathcal{A}^\bullet, \mathcal{O}_X)
$$
also follows in a straightforward manner from this.
(We observe that our choice of conventions regarding
Koszul complexes does intervene in the definition
of the isomorphism
$R\SheafHom_{\mathcal{O}_X}(\mathcal{O}_X, \mathcal{O}_Y) =
\wedge^r \mathcal{N}[-r]$.)
\end{proof}